\documentclass[12pt,a4paper]{article}
\usepackage{graphicx}
\usepackage{amsmath}
\usepackage{hyperref}
\usepackage{geometry}
\geometry{margin=2.5cm}

\title{\textbf{CDFD Runtime Paper 04: Engine, CLI, and Reproducible Execution}}
\author{Steve Bico Mujjabi, MD\\
\small Independent Researcher\\
\small ORCID: \href{https://orcid.org/0009-0001-0556-5516}{0009-0001-0556-5516}}
\date{May 2026}

\begin{document}
\maketitle

\begin{abstract}
\noindent \textbf{Executive synthesis.} The CDFD Runtime is the executable layer
of the CDFD program. This paper documents the current architecture: state
arrays in \texttt{engine/state.py}, physics and domain updates under
\texttt{engine/}, domain adapters under \texttt{domains/}, CDFL under
\texttt{dsl/}, shared command logic under \texttt{runtime/runner.py}, and the
CLI entrypoint \texttt{cdfd.py}. The CLI is the first serious interface because
it supports reproducible simulations, paper diagnostics, validation, export,
and automation before the web application is expanded.
\end{abstract}

\begin{figure}[ht]
\centering
\fbox{\begin{minipage}{0.92\textwidth}
\centering
\textbf{Runtime evidence path}\\[0.4em]
Prior CDFD mathematics $\longrightarrow$ CDFL/runtime state $\longrightarrow$ bounded output $\longrightarrow$ falsification target.
\end{minipage}}
\caption{Conceptual schematic for the runtime papers. The engine translates the mathematical frame from the prior CDFD papers into executable CDFL/runtime diagnostics; candidate outputs remain tied to explicit tests before any stronger claim is made.}
\end{figure}

\section{Prior CDFD Basis}
This manuscript does not rederive the mathematical core of CDFD. It uses the
Mujjabi Adaptive Operating Ratio and the constrained-vacuum notation introduced
in Part I \cite{MujjabiI2026}, the torus-knot hierarchy and boundary-mode work
from the Part I sequence \cite{MujjabiVI2026,MujjabiIX2026}, the public balance
law developed in the Part I capstone \cite{MujjabiX2026}, the Part I transport
tests \cite{MujjabiXII2026}, and the Life Number and tri-regime biology bridge
from Part II \cite{MujjabiPartII2026}. The runtime papers therefore act as an
execution layer: they keep the mathematics connected to commands, outputs,
falsification tests, and domain-specific limits.


\section{Architecture Principle}
The core engine must not depend on the web app. The CLI and web app call
the same backend, and both preserve a strict separation between:
\begin{itemize}
    \item numerical state evolution;
    \item domain interpretation;
    \item command presentation;
    \item output serialization;
    \item validation and finite-value checks.
\end{itemize}

\section{State and Kernel}
\texttt{engine/state.py} stores $\Phi$, $C$, $S$, $M_s$, $\Psi_s$, and
domain-specific metadata. \texttt{engine/kernel.py} orchestrates updates and
isolates domain failures so one module does not halt an entire run.

The kernel is deliberately conservative. It returns structured state, active
constraints, status flags, and metadata rather than hidden side effects.
Every domain module is optional; the physics step remains the common substrate.

\section{Shared Runtime Utilities}
The upgraded runtime exposes \texttt{runtime/diagnostics.py}. This module
centralizes tools that were previously scattered through release scripts:
\begin{itemize}
    \item strict JSON cleaning and non-finite audits;
    \item scalar and array forms of $\Psi_s$;
    \item the Life Number $\Lambda$;
    \item bounded adaptive updates for toy diagnostics;
    \item state summaries and result envelopes with provenance.
\end{itemize}

\section{CLI-First Interface}
The CLI entrypoint is \texttt{cdfd.py}. From the runtime root:
\begin{verbatim}
python cdfd.py info
python cdfd.py domains
python cdfd.py demo physics --steps 4 --nx 8 --ny 8 --json
python cdfd.py validate examples/heat_flow.cdfl
python cdfd.py run examples/heat_flow.cdfl --out outputs/heat_flow_run.json
python cdfd.py export outputs/heat_flow_run.json --out outputs/heat_flow_export.json
python cdfd.py auth --api-key-file app_key.txt --json
\end{verbatim}

The CLI calls \texttt{runtime/runner.py}. This makes the command backend reusable
for a future web app without duplicating model logic.

\section{Application Credentials}
The final CLI surface includes an app-auth boundary. A calling product can pass
an API key directly, through a key file, or through an environment variable.
The runtime never stores the raw key in output. It records only whether the key
matched the configured allowlist and, when a key is supplied, a short redacted
fingerprint for audit. This gives VOS and future apps a practical gate while
keeping LLM credentials and user-workflow policy out of the engine.

\section{Reproducible Output Contract}
Every serious command returns:
\begin{itemize}
    \item \texttt{status}: whether the command succeeded;
    \item \texttt{kind}: result type;
    \item \texttt{provenance}: runtime name, CDFL language tag, timestamp,
    command, Python version, and platform;
    \item \texttt{finite\_audit}: paths to non-finite values if any;
    \item \texttt{payload}: the domain, CDFL, or export result.
\end{itemize}
This is essential for paper diagnostics. A figure or discovery claim without
parameters, command, and output path is not part of the canonical runtime
evidence chain.

\section{Experiment Integration}
The experiments directory is not a separate mythology layer. It is the stress
test surface for the runtime:
\begin{itemize}
    \item \texttt{experiments/outputs/discovery\_vacuum\_hysteresis.h5};
    \item \texttt{experiments/outputs/discovery\_knot\_n7.h5};
    \item \texttt{experiments/outputs/discovery\_ool\_phase\_results.h5};
    \item \texttt{experiments/outputs/frontier\_sweep.h5};
    \item \texttt{experiments/outputs/gigamarathon\_1000\_results.h5}.
\end{itemize}
The report \texttt{experiments/reports/FINAL\_RELEASE\_EXPERIMENT\_SELECTION.md}
selects which outputs are suitable for public mention and how cautiously they
are framed.

\section{Current Verification}
The CLI/runtime upgrade pass was checked with:
\begin{verbatim}
PYTHONPATH=. /home/bampita/Projects/CDFD/.venv/bin/python \
  -m pytest -q tests/test_cli_runtime.py tests/test_new_engine_upgrades.py
\end{verbatim}
The recorded result was seven passing tests. This is not full validation of the
CDFD theory; it is evidence that the command backend, CDFL example, diagnostics
helpers, and selected native engine upgrades execute coherently.

\section{Conclusion}
The CDFD Runtime is becoming a reproducible scientific instrument rather than
only a manuscript companion. The CLI is the correct first interface because it
supports exact commands, saved outputs, test automation, and falsifiable
diagnostics before visual presentation is layered on top.
\begin{thebibliography}{9}
\bibitem{MujjabiI2026}
Steve Bico Mujjabi, MD. \emph{Vortex Stability in a Constrained Transport Vacuum and the Origin of the Fine-Structure Constant}. CDFD Part I Paper I, preprint, 2026.

\bibitem{MujjabiVI2026}
Steve Bico Mujjabi, MD. \emph{The Universal Torus Knot Hierarchy: $Z_n$ Symmetry, the Cinquefoil Family, and the General Structure of CDFT Particle Families}. CDFD Part I Paper VI, preprint, 2026.

\bibitem{MujjabiIX2026}
Steve Bico Mujjabi, MD. \emph{Even-$n$ Torus Knots in CDFT: The Exclusion of $T(2,2)$, the Extension to $T(2,2m)$, and the Anti-Phase Mass Scaling Law}. CDFD Part I Paper IX, preprint, 2026.

\bibitem{MujjabiX2026}
Steve Bico Mujjabi, MD. \emph{The Vacuum Equation of State: Deriving Torus Knot Self-Consistency from the Public CDFD Balance Law $\Psi = \Phi/C$}. CDFD Part I Paper X, preprint, 2026.

\bibitem{MujjabiXII2026}
Steve Bico Mujjabi, MD. \emph{Friction, Superconductivity, Plasma States, and Transport Mysteries: Observable Tests of Adaptive Constraint}. CDFD Part I Paper XII, preprint, 2026.

\bibitem{MujjabiPartII2026}
Steve Bico Mujjabi, MD. \emph{CDFD Part II: Origins of Life and Tri-Regime Bioenergetics}. Public release archive, 2026, \url{https://doi.org/10.5281/zenodo.20250821}.
\end{thebibliography}
\end{document}
